\section{Niepewność i ograniczenia oceny}
\label{sec:niepewnosc}
Rozstęp próbek opisuje zmienność zarejestrowanych wyników,
ale nie stanowi pełnej niepewności pomiaru. Przy ocenie wyniku
po korekcji tła należy uwzględnić rozrzut obu serii, przyrząd,
wzorcowanie, ustawienie mikrofonu, warunki propagacji i sposób
wyznaczenia tła. Instrukcja laboratoryjna nie podaje szczegółowego
algorytmu budżetu niepewności \autocite{piechowicz-konspekt}.

Dla modelu korekcji~\eqref{eq:korekcja-tla} współczynniki
wrażliwości na poziom całkowity i tło wynoszą
\begin{equation}
 c_m=\frac{10^{0{,}1L_{A\mathrm{sr}}}}
 {10^{0{,}1L_{A\mathrm{sr}}}-10^{0{,}1L_{At}}},\qquad
 c_t=-\frac{10^{0{,}1L_{At}}}
 {10^{0{,}1L_{A\mathrm{sr}}}-10^{0{,}1L_{At}}}.
\end{equation}
Wartości niepewności składowych i ich zależności trzeba ustalić
na podstawie dokumentacji zestawu i zastosowanej procedury.
Nie można bez uzasadnienia uznać wszystkich składników pomiaru
przy stacji i tła za niezależne, ponieważ użyto tego samego zestawu.

Nie wyznaczono liczbowej niepewności rozszerzonej: nie dysponowano
świadectwem wzorcowania ani wynikami sprawdzenia toru w dniu
pomiaru, a dokumentacja meteorologiczna i ocena porównywalności
punktu tła są niepełne. Przyjęcie arbitralnego błędu miernika lub
zerowej niepewności kalibracji nie byłoby uzasadnione.
Zwiększenie liczby cyfr wyniku nie usuwa tych ograniczeń.
